%============================================================ 선언문 (Elsevier 필수)
%  Elsevier 투고는 본문 뒤·참고문헌 앞에 아래 네 가지를 요구한다.
%    ① CRediT authorship contribution statement
%    ② Declaration of competing interest
%    ③ Declaration of generative AI and AI-assisted technologies in the writing process
%    ④ Data availability
%  ★ 절 제목은 Elsevier 가 정한 영문 문구 그대로여야 한다 — 제작 시스템이 제목으로
%    블록을 인식한다. 한글로 옮겨 적으면 규약 밖이다.
%
%  ★ ①은 저자만 아는 사실이라 초안으로 둔다. 붉은 표시가 남아 있으면 투고하지 말 것.
%    ②~④는 이 저장소에서 확인 가능한 사실로 채웠다.

\section*{CRediT authorship contribution statement}
\addcontentsline{toc}{section}{CRediT authorship contribution statement}

\textcolor{red}{\textbf{[투고 전 저자 확인 필요]}} 아래는 CRediT 분류에 따른 초안이며,
각 저자의 실제 기여로 교체하여야 한다.

\noindent
\textbf{윤영준}: 개념 설계, 방법론, 소프트웨어, 검증, 형식 분석, 조사, 자료 큐레이션,
초고 작성. \textbf{방현태}: 소프트웨어, 검증, 시각화. \textbf{전유라}: 방법론, 조사,
검토 및 편집. \textbf{김덕환}: 검토 및 편집, 감독. \textbf{윤원근}: 개념 설계, 감독,
연구 관리, 재원 획득.

\section*{Declaration of competing interest}

\textcolor{red}{\textbf{[투고 전 저자 확인 필요]}}
저자들은 본 논문에 보고된 연구에 영향을 미칠 수 있는 경쟁적 재정적 이해관계나 개인적 관계가
없음을 선언한다.

\section*{Declaration of generative AI and AI-assisted technologies in the writing process}

원고 작성 과정에서 저자들은 문헌 정리, 초고 문장 다듬기, \LaTeX{} 조판 및 도판 생성
스크립트 작성에 생성형 인공지능 보조 도구를 사용하였다. 모든 실험 설계, 데이터 수집,
수치 산출과 해석은 저자들이 수행하였으며, 저자들은 원고 전체를 검토하고 그 내용에 대해
전적인 책임을 진다.

\section*{Data availability}

본 연구의 결과를 재현하는 데 필요한 자료는 다음과 같이 구성된다.

\begin{itemize}[leftmargin=*, itemsep=2pt, topsep=3pt]
\item \textbf{원자료}: 720 에피소드의 에피소드별 지표(\texttt{episodes.csv}) 및 계획 호출
      계측값(\texttt{llm\_metrics.csv}). 본문의 모든 표와 그림은 이 두 파일에서
      기계적으로 생성되며(\texttt{tools/make\_paper\_tables.py},
      \texttt{boatattack\_sim/eval/paper\_figs.py}), 손으로 옮겨 적은 수치는 없다.
\item \textbf{정책 체크포인트}: 기동 계층의 가중치와 그 학습 설정(\texttt{u-net\_map.pt}).
      설정은 체크포인트 안에 함께 저장되어 있어 별도 설정 파일 없이 복원된다.
\item \textbf{시뮬레이터}: 환경, 관측 래스터화, 평가 하네스 전체.
\end{itemize}

\noindent
\textcolor{red}{\textbf{[투고 전 공개 URL 기입 필요]}} 위 자료는 논문 게재 시
공개 저장소를 통해 배포한다.

\vspace{3pt}
\noindent
\textbf{보존하지 않은 것.} 언어모델이 생성한 판단 근거(\texttt{rationale})의 \emph{원문}은
저장하지 않았고, 길이만 계측하였다. 따라서 근거 문장 자체에 대한 사후 분석은 본 자료로
재현할 수 없다(\S\ref{sec:limitations}). 또한 상용 API로 제공되는 폐쇄형 지휘관 모델은
버전이 갱신되거나 서비스가 종료될 수 있어, 그 조건의 수치가 시점에 따라 달라질 수 있다.
언어모델을 쓰지 않는 조건(\texttt{heur\_heur}, \texttt{heur\_unet})은 이 위험 없이
재현 가능하다.

